\clearpage
\fancyhead[L]{\small GERLING.}
\hypertarget{native-547-lower}{}\label{native-547-lower}
\begin{center}\small Göttingische gelehrte Anzeigen. Stück 190. S. 1889..1891. 28. November 1812.\\[3pt]\rule{.72\linewidth}{.3pt}\end{center}

\emph{Methodi proiectionis orthographicae usum ad calculos parallacticos facilitandos explicavit, simulque eclipsin solarem die VII. Sept. \textup{1820} apparituram hoc modo tractatam mappaque geographica illustratam tamquam exemplum proposuit} \textsc{Chr. Lud. Gerling}. 46 Seiten in Quart. Göttingen. Gedruckt bei \textsc{I. C. Baier}.

Um die Umstände einer Sternbedeckung oder Sonnenfinsternis für alle Punkte der Erde zu bestimmen, an denen sie sichtbar ist, verwenden Astronomen bekanntlich die orthografische Projektion auf eine Ebene, auf der die Gerade vom Mittelpunkt des bedeckten oder verfinsterten Gestirns senkrecht steht. Dieses Verfahren vermittelt eine leicht verständliche und klare Vorstellung von den verschiedenen Punkten im Raum, mit denen man es dabei zu tun hat. Es eröffnet damit einen bequemen Weg zur Lösung der auftretenden Aufgaben, ob man sich nun der Zeichnung oder der Rechnung, des schärfsten aller Instrumente, bedienen will. Die vorliegende kleine Abhandlung, die \SourceLeaf{548}ihr Verfasser --- früher unser gelehrter Mitbürger, heute Mathematiklehrer am königlichen Lyzeum in Cassel --- der hiesigen Philosophischen Fakultät anlässlich seiner Promotion als Probeschrift einreichte, bietet eine verständliche Anleitung zu dieser Methode und zeigt gründliche Kenntnisse. Die Herleitungen sind fast durchweg analytisch ausgeführt, sodass der Verfasser es nicht einmal für nötig hielt, eine Figur beizufügen. Dass er darauf verzichten konnte, ohne die Deutlichkeit zu beeinträchtigen, sodass jeder auch nur einigermaßen Geübte seinen Schlüssen leicht folgen kann, beweist die innere Klarheit, die der analytischen Methode bei geometrischen Aufgaben eigen ist, wenn sie zweckmäßig angewandt wird.

Die Sonnenfinsternis vom 7. September 1820, auf die der Verfasser die Methode beispielhaft anwendet, wird für unsere Gegend Europas auf lange Zeit eine der bemerkenswertesten sein. Die Orte, für die sie zentral ist, liegen auf einer Linie, die im nördlichsten Teil Amerikas beginnt, nördlich über Grönland und Island hinweg durch die Nordsee und in Richtung von Emden bis Triest durch Deutschland verläuft. Sie zieht sich dann der Länge nach durch das Adriatische Meer, über Morea und die nordlichste Spitze des Roten Meeres und endet bei Sonnenuntergang in Arabien. Auf dieser ganzen Linie und in einer ziemlich breiten Zone oberhalb und unterhalb davon ist die Finsternis ringförmig, auch bei uns in Göttingen. Dort wird der Ring nach der Rechnung des Verfassers von 2 Uhr 40 Min. 41 Sek. bis 2 Uhr 46 Min. 0 Sek. wahrer Zeit dauern. Auch für Bremen, Seeberg, Berlin, Wien und Mannheim hat der Verfasser hier die wichtigsten Umstände der Finsternis mitgeteilt. Für eine noch größere Zahl von Orten verspricht er, sie an anderer Stelle nachzutragen.
